% Part: lambda-calculus
% Chapter: introduction
% Section: free-variables

\documentclass[../../../include/open-logic-section]{subfiles}

\begin{document}

\olfileid{lam}{syn}{fv}
\olsection{মুক্ত চলরাশি}

ল্যাম্বডা ক্যালকুলাস অপেক্ষক নিয়ে আলোচনা করে, আর ল্যাম্বডা-বিমূর্তনের
মাধ্যমেই অপেক্ষকগুলি গড়ে ওঠে। স্বজ্ঞাতভাবে, $\lambd[x][M]$ হলো এমন
অপেক্ষক, যার আর্গুমেন্টকে~$x$-এ নির্ধারণ করলে মানগুলি~$M$ দিয়ে দেওয়া
হয়। কিন্তু~$M$-এ~$x$-এর প্রতিটি সংঘটন প্রাসঙ্গিক নয়: $M$-এর মধ্যে
আরেকটি বিমূর্তন $\lambd[x][N]$ থাকলে~$N$-এ~$x$-এর সংঘটনগুলি
$\lambd[x][N]$-এর জন্য প্রাসঙ্গিক, $\lambd[x][M]$-এর জন্য নয়। তাই
$\lambd[x][M]$-এর ভিতরের কোনো ল্যাম্বডা বিমূর্তন $\lambd[x]$, $M$-এ
$x$-এর যেসব সংঘটন অন্য কোনো ল্যাম্বডা বিমূর্তন ইতিমধ্যে আবদ্ধ করেনি—অর্থাৎ
$M$-এ~$x$-এর \emph{মুক্ত} সংঘটনগুলি—\emph{আবদ্ধ} করে।

\begin{defn}[পরিসর]
কোনো পদ~$N$-এর ভিতরে $\lambd[x][M]$ ঘটলে~$M$-এর সংশ্লিষ্ট সংঘটনটিই
$\lambd[x]$-এর \emph{পরিসর}।
% BN-SRC-275 TARGET-CORRECTION-BEGIN
\par\smallskip\noindent\textnormal{\small\textbf{উৎস-সংশোধন (BN-SRC-275):} হিমায়িত সংজ্ঞাটি~$\lambd[x][M]$-এর পরিসর হিসেবে ভিতরের দেহ~$M$-এর বদলে সমগ্র পরিবেষ্টক পদ~$N$-এর সংশ্লিষ্ট সংঘটনকে বলেছে। বাংলা পাঠে বিমূর্তনের গঠন ও পরের উদাহরণগুলি অনুসারে~$M$ পুনরুদ্ধার করা হয়েছে।} % BN-SRC-275 TARGET-CORRECTION-END
\end{defn}


\begin{defn}[মুক্ত ও বদ্ধ সংঘটন]
  কোনো পদ~$M$-এ চলরাশি~$x$-এর কোনো সংঘটন একটি $\lambd[x]$-এর পরিসরে
  না থাকলে সেটি \emph{মুক্ত}, অন্যথায় \emph{বদ্ধ}। $\lambd[x][M]$-এ
  কোনো চলরাশি~$x$-এর সংঘটন প্রারম্ভিক $\lambd[x]$ দ্বারা আবদ্ধ হয় যদি
  এবং কেবল যদি~$M$-এ~$x$-এর সেই সংঘটনটি মুক্ত হয়।
\end{defn}

\begin{ex}
  $\lambd[x][x y]$-এ $x$ ও~$y$ উভয়ই $\lambd[x]$-এর পরিসরে, তাই~$x$
  ওই $\lambd[x]$ দ্বারা আবদ্ধ। $y$ কোনো $\lambd[y]$-এর পরিসরে নেই বলে
  সেটি মুক্ত। $\lambd[x][x x]$-এ $x$-এর উভয় সংঘটনই $\lambd[x]$ দ্বারা
  আবদ্ধ, কারণ $xx$-তে দুটিই মুক্ত। $((\lambd[x][xx])x)$-এ~$x$-এর শেষ
  সংঘটনটি মুক্ত, কারণ সেটি কোনো $\lambd[x]$-এর পরিসরে নেই।
  $\lambd[x][(\lambd[x][x])x]$-এ প্রথম $\lambd[x]$-এর পরিসর হলো
  $(\lambd[x][x])x$, আর দ্বিতীয় $\lambd[x]$-এর পরিসর হলো~$x$-এর শেষের
  আগের সংঘটনটি। $(\lambd[x][x])x$-তে $x$-এর শেষ সংঘটনটি মুক্ত এবং শেষের
  আগেরটি বদ্ধ। তাই $\lambd[x][(\lambd[x][x])x]$-এ~$x$-এর শেষের আগের
  সংঘটনটি দ্বিতীয় $\lambd[x]$ দ্বারা এবং শেষটি প্রথম $\lambd[x]$ দ্বারা
  আবদ্ধ।
\end{ex}

কোনো পদ~$P$-এর সব চলরাশি-সংঘটন পরীক্ষা করে আমরা তার মুক্ত চলরাশিগুলির
সেট পাই। এই সেটটিকে $\FV{P}$ দিয়ে বোঝানো হয়; তার স্বাভাবিক সংজ্ঞাটি
নিচে দেওয়া হলো:

\begin{defn}[কোনো পদের মুক্ত চলরাশি] \ollabel{def:fv}
  কোনো পদের \emph{মুক্ত চলরাশিসমূহের} সেটটি নিচের নিয়মে আরোহীভাবে সংজ্ঞায়িত:
  \begin{enumerate}
    \item $\FV{x} = \{x\}$ \ollabel{def:fv1}
    \item $\FV{\lambd[x][N]} = \FV{N} \setminus \{x\}$    \ollabel{def:fv2}
    \item $\FV{PQ} = \FV{P} \cup \FV{Q}$ \ollabel{def:fv3}
  \end{enumerate}
\end{defn}

\begin{prob}
  \begin{enumerate}
  \item এই পদটিতে $\lambd[g]$ এবং দুটি $\lambd[x]$-এর পরিসর শনাক্ত করো:
    $\lambd[g][(\lambd[x][g (x x)]) \lambd[x][g (x x)]]$।
  \item $\lambd[g][(\lambd[x][g (x x)]) \lambd[x][g (x x)]]$-এ
    চলরাশির সব সংঘটন কি বদ্ধ? প্রতিটি কোন বিমূর্তন দ্বারা আবদ্ধ?
    \item $\FV{\lambd[x][(\lambd[y][(\lambd[z][x y]) z]) y]}$ নির্ণয় করো।
  \end{enumerate}
\end{prob}

\begin{explain}
মুক্ত চলরাশি বাইরের জগতের (\emph{পরিবেশের}) প্রতি একটি নির্দেশের মতো;
তাই মুক্ত চলরাশিযুক্ত কোনো পদকে আংশিকভাবে নির্দিষ্ট পদ হিসেবে দেখা যায়,
কারণ তার আচরণ নির্ভর করে আমরা পরিবেশটি কীভাবে স্থির করি তার উপর। যেমন,
$\lambd[x][f x]$ পদটি একটি আর্গুমেন্ট~$x$ গ্রহণ করে সেই আর্গুমেন্টের
উপর~$f$ প্রয়োগের ফল ফেরত দেয়; এখানে চলরাশি~$f$ মুক্ত। পদটির মান তার
পরিবেশের উপর, বিশেষত ওই পরিবেশে~$f$-এর মানের উপর, নির্ভরশীল।

এই পদের উপর বিমূর্তন প্রয়োগ করলে পাই $\lambd[f][\lambd[x][f
    x]]$। পদটি আর পরিবেশের চলরাশি~$f$-এর উপর নির্ভরশীল নয়, কারণ এখন এটি
এমন একটি অপেক্ষক নির্দেশ করে, যা দুটি আর্গুমেন্ট গ্রহণ করে এবং প্রথমটি
দ্বিতীয়টির উপর প্রয়োগ করার ফল ফেরত দেয়। পরিবেশে~$f$ বদলালে এই পদের
আচরণে কোনো প্রভাব পড়বে না; পদটি পরিবেশের~$f$-এর মান নয়, আর্গুমেন্ট
হিসেবে যা দেওয়া হবে কেবল সেটিই ব্যবহার করবে।
\end{explain}

\begin{defn}[বদ্ধ পদ, সমাবেশক]
  মুক্ত চলরাশিহীন পদকে \emph{বদ্ধ পদ}, অথবা \emph{সমাবেশক}, বলা হয়।
\end{defn}

\begin{lem}\ollabel{lem:fv}
  \begin{enumerate}
    \item \ollabel{lem:fv-abs} $y \neq x$ হলে $y \in
      \FV{\lambd[x][N]}$ যদি এবং কেবল যদি $y \in \FV{N}$।
    \item \ollabel{lem:fv-app} $y \in \FV{PQ}$ যদি এবং কেবল যদি
      $y \in \FV{P}$ অথবা $y \in \FV{Q}$।
    \end{enumerate}
\end{lem}

\begin{proof}
  অনুশীলনী।
\end{proof}

\begin{prob}
  \olref[lam][syn][fv]{lem:fv} প্রমাণ করো।
\end{prob}

\end{document}
